% !PRINTED 246
\weberpage{246}

If one of the two quantities $a_1$, $a_2$ vanishes, then, since in that case not all the roots can be of the same sign, we have one or two positive roots.

If $a_3 = 0$, the discriminant reduces to
\[
a_2^2(a_1^2 - 4a_2),
\]
and if this is positive, the cubic equation has one vanishing root and two further real roots. These are positive if
\[
a_1 < 0, \quad a_2 > 0,
\]
negative if
\[
a_1 > 0, \quad a_2 > 0,
\]
and one of them is positive and one negative if
\[
a_1 \gtreqless 0, \quad a_2 < 0
\]
holds.

\begin{center}
\S.\,78.

\gesperrt{Discussion of the biquadratic equation.}
\end{center}

For equations of the fourth and fifth degree there exist either no imaginary roots, or two, or four. If the discriminant is negative, there are two imaginary roots. With a positive discriminant, either four imaginary roots or none may be present. To distinguish these two cases will be our task in what follows. But first we turn to an elementary consideration of the biquadratic equation, which for the sake of simplicity we shall assume in the abbreviated form
\weq{1}{x^4 + ax^2 + bx + c = 0}
from which we can easily return to the general form.

If we denote the roots by $\alpha$, $\beta$, $\gamma$, $\delta$ and set, as in \S.\,47,
\weq{2}{\begin{aligned} \alpha - \beta &= v+w, & \gamma - \delta &= v-w,\\ \alpha - \gamma &= w+u, & \delta - \beta &= w-u,\\ \alpha - \delta &= u+v, & \beta - \gamma &= u-v, \end{aligned}}
then $u^2$, $v^2$, $w^2$ are the roots of the cubic resolvent
\weq{3}{y^3 + 2ay^2 + (a^2-4c)y - b^2 = 0,}
and the discriminant $D$ of this cubic equation, which by
%CONT>>
